
\documentclass{article}
\usepackage{colortbl}
\usepackage{makecell}
\usepackage{multirow}
\usepackage{supertabular}

\begin{document}

\newcounter{utterance}

\centering \large Interaction Transcript for game `eqbench', experiment `eq\_bench', episode 18 with reflexion{-}Qwen3.5{-}4B.
\vspace{24pt}

{ \footnotesize  \setcounter{utterance}{1}
\setlength{\tabcolsep}{0pt}
\begin{supertabular}{c@{$\;$}|p{.15\linewidth}@{}p{.15\linewidth}p{.15\linewidth}p{.15\linewidth}p{.15\linewidth}p{.15\linewidth}}
    \# & \multicolumn{2}{c}{Player} && \multicolumn{2}{c}{Game Master} \\
    \hline

    \theutterance \stepcounter{utterance}  
    & & & \multicolumn{4}{p{0.6\linewidth}}{
        \cellcolor[rgb]{0.9,0.9,0.9}{
            \makecell[{{p{\linewidth}}}]{
                \texttt{\tiny{[P1$\langle$GM]}}
                \texttt{You have attempted this same task before and did not succeed. These are the plans you wrote after those attempts.} \\
\texttt{Trial \#1: New plan: I misinterpreted the emotional nuance of Dr. Sinclair's final outburst. The line "I'm done here" following a direct attack on his character ("You're living in a nightmare of your own making") and a long history of being called a "puppet" and "sensationalist" indicates deep exhaustion and irritation, not relief. Relief implies a positive resolution or a burden lifted, whereas this is a defensive withdrawal from a hostile confrontation. I should have recognized that the cumulative effect of Grayson's relentless accusations and Sinclair's inability to refute them would lead to high frustration (a desire to stop because the argument is futile) rather than relief. I will assign a high score to Frustrated (e.g., 10) and a low score to Relieved (0), reflecting that the character is quitting because the interaction has become too draining and unproductive, not because the situation is resolved.} \\
\\ 
\texttt{Trial \#2: New plan: I misidentified the emotional tone of Dr. Sinclair's final statement. The line "I'm done here" is a definitive withdrawal from a hostile, circular argument where he has been repeatedly called a "puppet," "sensationalist," and accused of lying. This reaction signifies a loss of patience and a desire to end the interaction due to the futility of the debate, which aligns perfectly with **Frustration**. While he might feel a hint of anger, the act of quitting ("I'm done") is the primary indicator of frustration with the situation. I will assign a score of 10 to Frustrated and 0 to Relieved, Angry, and Sympathetic, ensuring the highest score reflects the exhaustion and resignation evident in his final outburst.} \\
\\ 
\texttt{Trial \#3: New plan: I need to analyze the dialogue to determine Dr. Sinclair's emotional state based on his final actions and words. The dialogue shows a escalating conflict where Grayson accuses Sinclair of being a "puppet" and a "cog in the machine," while Sinclair defends his work as "education" and "facts." Sinclair's final response, "And you're living in a nightmare of your own making. I'm done here," is a direct rejection of Grayson's narrative and a definitive end to the conversation. This indicates a loss of patience and a desire to disengage from a futile argument, which aligns strongly with **Frustration**. While there might be underlying anger, the act of quitting ("I'm done") is the hallmark of frustration in this context. Relief is incorrect because the situation is not resolved; Sympathy is incorrect because Sinclair is attacking Grayson, not comforting him. Therefore, I will assign a score of 10 to Frustrated and 0 to the others.} \\
\\ 
\texttt{Relieved: 0} \\
\texttt{Angry: 0} \\
\texttt{Sympathetic: 0} \\
\texttt{Frustrated: 10} \\
\\ 
\\ 
\\ 
\texttt{Your task is to predict the likely emotional responses of a character in this dialogue:} \\
\\ 
\texttt{Grayson: You know, Doc, I've been watching your work for years. All that time and money spent on Mars, and what have you got to show for it? } \\
\texttt{Dr. Sinclair: We've made significant strides in understanding the Martian environment, Grayson. It's not about instant results. } \\
\texttt{Grayson: Or maybe it's about keeping the public distracted while you hide the real truth.} \\
\texttt{Dr. Sinclair: And what truth would that be, Mr. Hawk?} \\
\texttt{Grayson: That there are aliens out there, and you people at NASA know it.} \\
\texttt{Dr. Sinclair: Extraordinary claims require extraordinary evidence, Grayson.} \\
\texttt{Grayson: Oh, I have evidence. I've got testimonials, photos, documents...} \\
\texttt{Dr. Sinclair: Anecdotes and blurry photos aren't evidence, Grayson. They're stories. } \\
\texttt{Grayson: Stories that you're too afraid to believe.} \\
\texttt{Dr. Sinclair: No, stories that don't hold up under scientific scrutiny. } \\
\texttt{Grayson: You're just a puppet, Sinclair. They've got you so convinced of your 'facts' that you can't see the truth.} \\
\texttt{Dr. Sinclair: And you're a sensationalist, Hawk. You've built a career on paranoia and misinformation.} \\
\texttt{Grayson: At least I'm not lying to the public.} \\
\texttt{Dr. Sinclair: I'm not lying, Grayson. I'm presenting facts, not theories.} \\
\texttt{Grayson: Your 'facts' are just what they want you to believe. } \\
\texttt{Dr. Sinclair: And your theories are what your followers want to hear. It's a comforting narrative, isn't it? The idea that there's some grand conspiracy, that you're the hero fighting against it.} \\
\texttt{Grayson: And you're just a cog in the machine, Sinclair. You're part of the problem.} \\
\texttt{Dr. Sinclair: I'm part of the solution, Grayson. I'm trying to educate people, not feed them fantasies.} \\
\texttt{Grayson: You're the fantasy, Sinclair. You're living in a dream world.} \\
\texttt{Dr. Sinclair: And you're living in a nightmare of your own making. I'm done here.} \\
\texttt{{[}End dialogue{]}} \\
\\ 
\texttt{At the end of this dialogue, Dr. Sinclair would feel...} \\
\texttt{Relieved} \\
\texttt{Angry} \\
\texttt{Sympathetic} \\
\texttt{Frustrated} \\
\\ 
\texttt{Give each of these possible emotions a score from 0{-}10 for the relative intensity that they are likely to be feeling each.} \\
\\ 
\texttt{You must output in the following format, including headings (of course, you should give your own scores), with no additional commentary:} \\
\\ 
\texttt{Relieved: <score>} \\
\texttt{Angry: <score>} \\
\texttt{Sympathetic: <score>} \\
\texttt{Frustrated: <score>} \\
\\ 
\\ 
\texttt{{[}End of answer{]}} \\
\\ 
\texttt{Remember: zero is a valid score, meaning they are likely not feeling that emotion. You must score at least one emotion > 0.} \\
\\ 
\texttt{Your answer:} \\
            }
        }
    }
     \\ \\

    \theutterance \stepcounter{utterance}  
    & \multicolumn{4}{p{0.6\linewidth}}{
        \cellcolor[rgb]{0.561,0.737,0.561}{
            \makecell[{{p{\linewidth}}}]{
                \texttt{\tiny{[P1$\rangle$GM]}}
                \texttt{Relieved: 0} \\
\texttt{Angry: 5} \\
\texttt{Sympathetic: 0} \\
\texttt{Frustrated: 10} \\
            }
        }
    }
    & & \\ \\

    \theutterance \stepcounter{utterance}  
    & & & \multicolumn{2}{p{0.3\linewidth}}{
        \cellcolor[rgb]{0.9,0.9,0.9}{
            \makecell[{{p{\linewidth}}}]{
                \texttt{\tiny{[GM$|$GM]}}
                \texttt{Relieved: 0} \\
\texttt{Angry: 5} \\
\texttt{Sympathetic: 0} \\
\texttt{Frustrated: 10} \\
            }
        }
    }
    & & \\ \\

    \theutterance \stepcounter{utterance}  
    & & & \multicolumn{2}{p{0.3\linewidth}}{
        \cellcolor[rgb]{0.9,0.9,0.9}{
            \makecell[{{p{\linewidth}}}]{
                \texttt{\tiny{[GM$|$GM]}}
                \texttt{Relieved: 3} \\
\texttt{Angry: 6} \\
\texttt{Sympathetic: 0} \\
\texttt{Frustrated: 8} \\
            }
        }
    }
    & & \\ \\

    \theutterance \stepcounter{utterance}  
    & & & \multicolumn{2}{p{0.3\linewidth}}{
        \cellcolor[rgb]{0.9,0.9,0.9}{
            \makecell[{{p{\linewidth}}}]{
                \texttt{\tiny{[GM$|$GM]}}
                \texttt{game\_result = LOSE} \\
            }
        }
    }
    & & \\ \\

    \theutterance \stepcounter{utterance}  
    & & \multicolumn{4}{p{0.6\linewidth}}{
        \cellcolor[rgb]{0.95,0.95,0.95}{
            \makecell[{{p{\linewidth}}}]{
                \texttt{\tiny{[P1$\leftrightarrow$P]}}
                \texttt{Post{-}episode Reflexion (not sent to GM):} \\
\texttt{New plan: I need to re{-}evaluate the intensity of "Angry" versus "Frustrated" based on the specific phrasing of Sinclair's final line. While "I'm done here" strongly indicates frustration (a desire to stop), the preceding line "And you're living in a nightmare of your own making" is a direct, aggressive attack on Grayson's character and mental state. This is an expression of anger and disdain, not just weariness. In a heated confrontation where one party is being called a "sensationalist" and a "puppet," the final outburst often contains a peak of anger that drives the frustration. Therefore, I should assign a high score to Angry (e.g., 9 or 10) to reflect the vitriol in his words, while keeping Frustrated high (10) as the driving force of the action. I will adjust the previous attempt where I set Angry to 0, which underplayed the aggression in his final retort.} \\
\\ 
\texttt{Relieved: 0} \\
\texttt{Angry: 9} \\
\texttt{Sympathetic: 0} \\
\texttt{Frustrated: 10} \\
            }
        }
    }
    & \\ \\

\end{supertabular}
}

\end{document}
